\documentclass[conference]{IEEEtran}
\IEEEoverridecommandlockouts

\usepackage{cite}
\usepackage{amsmath,amssymb,amsfonts,amsthm}
\usepackage{algorithm}
\usepackage{algorithmic}

\usepackage{graphicx}
\graphicspath{{figures/pics/}}
\usepackage{caption}
\usepackage{booktabs}
\usepackage{multirow}

\usepackage{textcomp}
\usepackage{xcolor}
\usepackage{url}
\usepackage{microtype}
\usepackage[hidelinks]{hyperref}

\newtheorem{definition}{Definition}
\newtheorem{proposition}{Proposition}
\newtheorem{remark}{Remark}
\captionsetup[figure]{font=footnotesize}

\newcommand{\safeincludegraphics}[2][]{%
  \IfFileExists{#2}{%
    \includegraphics[#1]{#2}%
  }{%
    \fbox{\parbox[c][0.17\textheight][c]{0.9\linewidth}{%
      \centering\small Figure placeholder:\\ \texttt{\detokenize{#2}}}}%
  }%
}

\newcommand{\figcap}[1]{\caption{#1}}

\def\BibTeX{{\rm B\kern-.05em{\sc i\kern-.025em b}\kern-.08em
    T\kern-.1667em\lower.7ex\hbox{E}\kern-.125emX}}

\begin{document}

\title{Deterministic Retrieval for FinTech:\\
Formalizing Stratified Knowledge Graphs\\
for Zero Hallucination Compliance}

\author{%
  \IEEEauthorblockN{Muhammad Adil Usmani}
  \IEEEauthorblockA{%
    \textit{Department of Computer Science}\\
    \textit{University of Central Punjab}\\
    Lahore, Pakistan\\
    \href{mailto:muhammadaadilusmani@gmail.com}{muhammadaadilusmani@gmail.com}
  }
}

\maketitle

\begin{abstract}
Financial compliance workflows require verifiable provenance, reproducible execution, and conservative failure behavior. Standard Retrieval Augmented Generation pipelines often retrieve broad semantically similar passages that increase context noise and permit unsupported synthesis. This paper presents a Lexical Structural GraphRAG architecture for SEC 10-K analysis that combines strict lexical anchoring with deterministic graph traversal. We formalize two principles. The first is Triplet Atomicity, which requires evidence to be represented and retrieved as typed subject predicate object facts rather than mixed prose chunks. The second is a Deterministic Path Requirement, which prohibits answer synthesis when no bounded structural path connects the query entities in the knowledge graph. On a 35 query benchmark spanning direct, implicit, comparative, and multi hop financial questions, the system achieves 100\% execution success, mean latency of about 18 seconds, Recall@K of 0.86, faithfulness of 0.97, and 0\% hallucination on unretrieved context through strict refusal. Results show that triplet level grounding with deterministic gating yields stronger audit readiness and comparative reasoning than vector only baselines in high risk compliance settings.
\end{abstract}

\begin{IEEEkeywords}
GraphRAG, Triplet Atomicity, FinTech Compliance, Deterministic Retrieval, SEC 10-K, Explainable AI, Knowledge Graphs, Hallucination Mitigation, Information Retrieval.
\end{IEEEkeywords}

\section{Introduction}
The central challenge in deploying large language models for financial operations is not fluency but verifiable reliability. In regulated environments, every generated statement must be traceable to source evidence and defensible under audit \cite{basel,gdpr,sr117,bcbs239,iso27001,pci40}. Conventional Retrieval Augmented Generation improves factual grounding relative to closed book generation, yet it still inherits a major weakness: semantic proximity is not the same as evidential sufficiency \cite{lewis2020rag,karpukhin2020dpr,colbert2020,splade2021,atlas2022}. When large text chunks are inserted into model context, adjacent irrelevant sentences create noise, and the model may interpolate unsupported claims in a way that is difficult to inspect after the fact.

This paper proposes a shift from paragraph retrieval to graph based symbolic retrieval. The proposed method combines three mechanisms. First, \textbf{Triplet Atomicity} stores evidence as explicit typed facts of the form subject, predicate, and object. Second, the \textbf{Deterministic Path Requirement} allows answer generation only when a bounded structural path exists between anchored query entities. Third, \textbf{Stratified Scoring} ranks candidate evidence by combining lexical intent with controlled centrality so high degree institutional nodes do not dominate retrieval neighborhoods.

The conservative refusal policy follows documented concerns about stochastic generation risk \cite{bender2021stochastic}. It also matches the compliance first requirements of SEC filing analysis \cite{sec10k,projectnotes}. In this setting, a false refusal can be escalated to a human analyst, while a hallucinated compliance claim can create regulatory, legal, and capital risk.

\subsection*{Contributions}
This paper makes five contributions.
\begin{enumerate}
  \item It defines Triplet Atomicity and places it within a provenance mapped financial knowledge graph.
  \item It formalizes a Deterministic Path Requirement with bounded hop traversal and stratified scoring for mega node control.
  \item It presents an end to end Lexical Structural GraphRAG pipeline for SEC 10-K analysis, including ingestion, semantic chunking, triplet extraction, canonicalization, and refusal gated synthesis.
  \item It evaluates the system on a 35 query benchmark with ablation, error taxonomy, token efficiency analysis, and latency analysis.
  \item It maps the architecture to SR 11-7, GDPR, and emerging high risk AI governance expectations.
\end{enumerate}

\section{Related Work}
\subsection{A Taxonomy of RAG Systems}
Current RAG systems can be grouped into three broad classes. \textit{Naive RAG} retrieves top $k$ chunks by vector similarity and forwards them directly to the generator \cite{lewis2020rag}. \textit{Modular RAG} adds reranking, query rewriting, filtering, and hybrid retrieval components \cite{karpukhin2020dpr,colbert2020,splade2021,hyde2023}. \textit{GraphRAG} introduces explicit entity and relation structure so retrieval can follow graph level paths rather than only semantic neighborhoods \cite{edge2024graphrag}. The present work belongs to the third class but differs by using deterministic lexical entry constraints, refusal gating before synthesis, and a scoring rule designed for financial mega node suppression.

These families emerge from transformer era representation learning and instruction tuned generation pipelines \cite{attention2017,bert2019,roberta2019,sbert2019,t52020,instructiontuning2022}. Our concern is not only answer quality but also the audit properties of the retrieval path that leads to each answer.

\subsection{Hallucination, Context Position, and Evidence Dilution}
Even strong language models can fail to use evidence that appears in the middle of long context windows. Liu et al. describe this as the ``Lost in the Middle'' effect \cite{liu2024lostmiddle}. Other work on truthfulness, retrieval faithfulness, and benchmark robustness shows that apparently strong performance can mask unsupported generation \cite{truthfulqa2022,ragas2023,hotpotqa2018,fever2018}. In compliance analysis this is especially costly because a single buried covenant or regulatory clause may determine whether a generated answer is valid.

Triplet Atomicity addresses this problem by replacing long mixed paragraphs with compact relation focused evidence units. Instead of asking the model to discover latent dependencies inside noisy prose, the system retrieves pre resolved typed facts with explicit provenance. This reduces effective context length and makes evidence boundaries inspectable.

\subsection{Knowledge Graphs for Financial Documents}
Prior work on financial knowledge graphs and EDGAR based graph construction shows the value of entity normalization, relation extraction, and graph analytics for risk monitoring \cite{edgargraph,finkgsurvey,opencyc2018,finbert2020}. However, most existing systems emphasize analytics dashboards or document intelligence rather than retrieval constrained generation. Our Lexical Structural hybrid differs in three ways: lexical anchoring controls node entry, deterministic path checks control synthesis permission, and refusal is treated as an auditable safety control rather than a retrieval failure.

\section{Architectural Structure}
The proposed pipeline operates in four phases: document ingestion and preprocessing, triplet extraction and graph construction, entity canonicalization, and lexical anchoring with deterministic traversal. Every phase preserves a provenance link so generated claims can be traced back to exact source chunks.

\begin{figure}[htbp]
  \centering
  \safeincludegraphics[width=0.94\linewidth]{image_d5e5de.png.png}
  \figcap{High level system schematic. The diagram summarizes the movement from query parsing to entity extraction, structural traversal, and grounded answer synthesis.}
  \label{fig4_overview_architecture}
\end{figure}

\begin{figure*}[!t]
  \centering
  \safeincludegraphics[width=0.86\textwidth]{plot_8_pipeline_comparison.png}
  \figcap{End to end pipeline view. Left, conventional vector centric retrieval injects prose chunks into the generator context without structural constraint. Right, the proposed Lexical Structural pipeline enforces lexical node anchoring, bounded traversal with degree caps, and explicit refusal before synthesis. Each generated claim is backed by a \texttt{source\_chunk} pointer.}
  \label{fig8_pipeline_main}
\end{figure*}

\subsection{Document Ingestion and Preprocessing}
SEC 10-K filings are fetched through an asynchronous ingestion routine with configurable retry logic, section tracking, and metadata normalization. HTML is parsed into semantically coherent text units, and each chunk stores a filing identifier, issuer identifier, section name, and page or offset metadata where available. This design supports efficient multi document processing while preserving traceability to the original filing structure.

\subsection{Chunking Strategy: Semantic vs. Structural}
Unlike fixed sliding windows, semantic chunking preserves topical boundaries in risk disclosures, management discussion, and notes to financial statements. Chunks are closed at topic shifts rather than a strict token cutoff. This reduces cross topic contamination before triplet construction and improves extraction precision relative to random chunking, as confirmed by the ablation in Section~V.

\begin{figure}[htbp]
  \centering
  \safeincludegraphics[width=0.95\linewidth]{plot_6_signal_to_noise.png}
  \figcap{Signal to noise comparison between paragraph retrieval and triplet atomic retrieval. Atomic evidence preserves a higher share of relevant tokens and lowers prompt overhead across query categories.}
  \label{fig6_signal_noise}
\end{figure}

\subsection{Triplet Extraction and Ontology}
Using \texttt{gpt-4o-mini} with low temperature for structured output stability, each chunk is transformed into a typed ontology with four node classes: \textit{Entity}, \textit{Obligation}, \textit{RiskFactor}, and \textit{Metric}. Edges encode typed predicates such as \textsc{reports}, \textsc{subject\_to}, \textsc{has\_risk}, and \textsc{discloses}. Every node stores a \texttt{source\_chunk} key so claims remain reverse auditable to the originating filing passage.

\begin{figure}[htbp]
  \centering
  \safeincludegraphics[width=0.96\linewidth]{plot_7_structured_knowledge_graph.png}
  \figcap{Structured financial knowledge graph topology used during deterministic traversal. Nodes represent canonical entities, obligations, risk factors, and metrics. Directed edges represent auditable typed predicates linked to source chunks.}
  \label{fig7_structured_kg}
\end{figure}

\subsection{Entity Resolution and Canonicalization}
Canonicalization applies three resolution layers in order. First, exact alias dictionaries collapse known variants such as ``JPM,'' ``JPMorgan,'' and ``JPMorgan Chase \& Co.'' into a single canonical key. Second, normalized string matching applies lowercasing, punctuation stripping, and legal suffix removal. Third, fuzzy overlap matching based on token level Jaccard similarity is used only when both deterministic layers fail. Ambiguous merges are rejected by default because a false merge can propagate incorrect facts across an entire graph neighborhood.

\subsection{Implementation: Neo4j and LLM Configuration}
The graph layer runs on Neo4j with composite indexes on canonical entity names, node labels, and document identifiers \cite{neo4jmanual,cypher2018}. Node properties include \texttt{canonical\_name}, \texttt{aliases}, \texttt{source\_chunk}, \texttt{document\_id}, \texttt{document\_type}, \texttt{section}, \texttt{page\_number}, and optional \texttt{vector\_embedding}. Edge properties include \texttt{predicate}, \texttt{confidence}, \texttt{pagerank}, and \texttt{timestamp}. The extraction model uses temperature $T=0.1$ and a strict JSON style output schema to reduce run to run variance \cite{gpt4report,instructiontuning2022}.

\begin{figure}[htbp]
  \centering
  \safeincludegraphics[width=0.96\linewidth]{plot_robust_architecture_plot_10.png}
  \figcap{Robust architecture view. The control path emphasizes ingestion discipline, canonicalization, deterministic traversal, provenance preservation, and refusal safe generation.}
  \label{fig10_robust_architecture}
\end{figure}

\section{Formalizing Deterministic Retrieval}
\subsection{Graph Definition and Provenance Mapping}
\begin{definition}[Provenance Knowledge Graph]
Let the financial knowledge base be a labeled property graph
\begin{equation}
  G = (V, E, L, \Phi),
\end{equation}
where $V$ is the node set, $E \subseteq V \times V$ is the directed edge set, $L$ is a typed label space over entity classes and relation predicates, and $\Phi : V \cup E \rightarrow \mathcal{C}$ maps every graph element to the source chunk set $\mathcal{C}$.
\end{definition}

\begin{definition}[Triplet Atomicity]
A triplet $t = (s,p,o)$ satisfies Triplet Atomicity if and only if
\begin{equation}
  \Phi(t) = c \in \mathcal{C}
\end{equation}
and the claim cannot be decomposed into sub claims with different provenance mappings. In practice this means each retrieved fact remains attached to one auditable evidence unit.
\end{definition}

\begin{remark}
Triplet Atomicity is stricter than generic triplet extraction. It prevents relation synthesis across chunk boundaries, which is a common source of unsupported inference.
\end{remark}

\subsection{Extraction Objective}
Given a source chunk $c \in \mathcal{C}$, the extraction model induces a candidate triplet set $T(c)$. The operational objective is to maximize factual extraction fidelity while minimizing unsupported relation injection:
\begin{equation}
  \min_{\theta} \; \mathbb{P}(\text{Hallucination} \mid c)
  \quad \text{subject to} \quad \Phi(t)=c, \; \forall t \in T(c).
\end{equation}
Although no custom model training is performed, this objective governs prompt design, schema constraints, and post extraction validation.

\subsection{Deterministic Path Requirement}
\begin{definition}[Valid Structural Path]
For anchored query entities $u,v \in V$, a valid structural path exists if
\begin{equation}
  \exists\; \pi(u,v) = \langle e_1, e_2, \ldots, e_k \rangle \subseteq E
  \quad \text{s.t.} \quad d(u,v)=k \leq \tau,
\end{equation}
where $\tau$ is the maximum hop bound and $d(u,v)$ is the shortest directed path length.
\end{definition}

\begin{proposition}[Refusal Correctness]
If no valid structural path exists for any anchored entity pair derived from a query, the system emits an explicit refusal rather than a speculative completion. Therefore every generated answer is backed by at least one retrievable provenance chain.
\end{proposition}

The refusal predicate is
\begin{equation}
  \mathcal{R}(q) =
  \begin{cases}
    \textsc{refusal}, & \text{if } A = \emptyset \text{ or no valid path exists},\\
    \textsc{answer}, & \text{otherwise},
  \end{cases}
\end{equation}
where $A$ is the anchor set resolved from the query.

\subsection{Mega Node Mitigation and Stratified Scoring}
High degree nodes that represent major institutions can dominate graph neighborhoods and bias retrieval toward frequent mentions rather than relevant evidence. To control this effect, candidate neighbors are first filtered by a degree cap $\deg(n) < \kappa$ and then ranked by
\begin{equation}
  S_{\mathrm{final}}(e)=W_{\mathrm{lexical}}(e)+\ln(\mathrm{PR}(e)+1),
\end{equation}
where $W_{\mathrm{lexical}}(e)$ is a normalized lexical overlap score and $\mathrm{PR}(e)$ is the PageRank score of the target node \cite{angles2018propertygraphs}. The logarithmic transform prevents centrality from overwhelming lexical relevance.

\begin{figure*}[!t]
  \centering
  \safeincludegraphics[width=0.98\textwidth]{plot_robust_architecture_plot_10.png}
  \figcap{Wide control view of the robust retrieval architecture. The system reduces hub dominance through lexical anchoring, degree limits, and path bounded evidence selection.}
  \label{fig10_robust_architecture_wide}
\end{figure*}

\subsection{Retrieval Algorithm}
\begin{algorithm}[htbp]
  \caption{Deterministic Lexical Structural Retrieval}
  \label{alg:deterministic_retrieval}
  \begin{algorithmic}[1]
    \REQUIRE Query $q$, graph $G=(V,E,L,\Phi)$, hop bound $\tau$, degree cap $\kappa$, top $k$
    \ENSURE Answer with source pointers, or explicit \textsc{refusal}
    \STATE Extract entity mentions $Q_E$ from $q$
    \STATE Resolve $Q_E$ to canonical anchor set $A \subseteq V$
    \IF{$A = \emptyset$}
      \STATE \textbf{return} \textsc{refusal}
    \ENDIF
    \STATE Build candidate subgraph with $\deg(n) < \kappa$
    \STATE Execute bounded traversal from each anchor with depth $\leq \tau$
    \IF{no valid path satisfies query constraints}
      \STATE \textbf{return} \textsc{refusal}
    \ENDIF
    \FOR{each candidate edge $e$ on retrieved paths}
      \STATE Compute $S_{\mathrm{final}}(e)$
    \ENDFOR
    \STATE Select top $k$ triplets ranked by $S_{\mathrm{final}}$
    \STATE Collect provenance mappings $\Phi(t)$ for all selected triplets
    \STATE Synthesize answer constrained strictly to retrieved triplets
    \STATE \textbf{return} \textsc{answer} with \texttt{source\_chunk} pointers
  \end{algorithmic}
\end{algorithm}

\noindent\textbf{Complexity.} Graph construction is $O(|V|+|E|)$. Bounded traversal from $|A|$ anchors with depth $\tau$ runs in $O(|A| \cdot b^{\tau})$, where $b$ is the average branching factor after degree capping. For small hop bounds used in compliance retrieval, this remains tractable in real time.

\subsection{Evaluation Metrics}
Retrieval quality is measured by contextual recall over expected evidence chunks \cite{msmarco2016,trec2020}:
\begin{equation}
  \mathrm{Recall@K} =
  \frac{|\mathcal{R}_K \cap \mathcal{E}|}{|\mathcal{E}|},
\end{equation}
where $\mathcal{R}_K$ is the set of top $K$ retrieved chunks and $\mathcal{E}$ is the set of annotated expected evidence chunks. Generation quality is assessed by faithfulness, hallucination rate, and refusal correctness.

\section{Performance, Ablation, and Failure Analysis}
\subsection{Benchmark Design and Query Categories}
The 35 query benchmark was constructed across four difficulty tiers. Direct queries require a single fact lookup. Implicit queries require lexical bridging between user vocabulary and canonical graph labels. Comparative queries require synchronized traversal from two entity anchors. Multi hop queries require valid paths across one or more intermediate nodes.

\begin{table}[htbp]
  \caption{Benchmark Query Category Breakdown}
  \label{tab:query_cats}
  \centering
  \begin{tabular}{lccc}
    \toprule
    Category & Count & Avg. Hops & Refusal Rate \\
    \midrule
    Direct & 12 & 1.0 & 8.3\% \\
    Implicit & 8 & 1.5 & 62.5\% \\
    Comparative & 9 & 2.1 & 44.4\% \\
    Multi hop & 6 & 2.8 & 33.3\% \\
    \midrule
    Overall & 35 & 1.8 & 42.8\% \\
    \bottomrule
  \end{tabular}
\end{table}

\subsection{Latency and Throughput}
Across all benchmark queries, the system achieved 100\% execution success with mean end to end latency of about 18 seconds per query. The dominant contributor was answer synthesis, with graph traversal and anchor resolution accounting for the remaining time. Observed traversal latency grew sub linearly with corpus size because of label and property indexes together with bounded hop search \cite{angles2018propertygraphs,robinson2015graphdb}.

\begin{figure}[htbp]
  \centering
  \safeincludegraphics[width=1.0\linewidth]{plot_fintech_latency_hallucination_tradeoff_plot_9.png}
  \figcap{Latency and hallucination trade off for financial retrieval pipelines. The proposed method accepts moderate latency overhead to reduce unsupported generation to zero on unretrieved context.}
  \label{fig_latency_tradeoff}
\end{figure}

\subsection{Ablation Study}
Table~\ref{tab:ablation} isolates the effect of the main architectural constraints. Removing PageRank reweighting increased hub dominance and reduced comparative precision. Replacing semantic chunking with random chunking degraded extraction quality and increased refusal due to fragmented evidence. Disabling refusal yielded the highest coverage but also the largest hallucination rate. These trends match prior findings that reranking, calibration, and retrieval structure materially affect downstream faithfulness \cite{monoT52020,rankgpt2023,selfrag2023}.

\begin{table}[htbp]
  \caption{Ablation Results on the 35 Query Benchmark}
  \label{tab:ablation}
  \centering
  \begin{tabular}{lccc}
    \toprule
    Configuration & Recall@K & Faithfulness & Hallucination \\
    \midrule
    Full system & 0.86 & 0.97 & 0.00 \\
    No PageRank term & 0.83 & 0.92 & 0.03 \\
    Random chunking & 0.74 & 0.90 & 0.04 \\
    No refusal gate & 0.87 & 0.81 & 0.14 \\
    \bottomrule
  \end{tabular}
\end{table}

\subsection{Precision vs. Hop Complexity}
Vector baselines improve apparent recall by expanding retrieved text volume, but precision drops as hop complexity increases because semantically nearby chunks introduce cross entity noise. The graph constrained method maintains higher precision by returning only path relevant atomic relations.

\begin{figure}[htbp]
  \centering
  \safeincludegraphics[width=0.8\linewidth]{fig1_precision_degradation.png}
  \captionsetup{font=scriptsize}
  \figcap{Precision degradation under increasing query complexity. Vector only retrieval degrades more sharply at higher hop complexity than lexical structural retrieval.}
  \label{fig1}
\end{figure}

\begin{figure}[htbp]
  \centering
  \safeincludegraphics[width=0.8\linewidth]{fig2_behavior_complexity.png}
  \figcap{System behavior across direct, comparative, and multi hop query categories. Refusal rises where lexical bridging or multi anchor traversal is harder to satisfy.}
  \label{fig2}
\end{figure}

\begin{figure}[htbp]
  \centering
  \safeincludegraphics[width=0.88\linewidth]{plot_5_cumulative_recall.png}
  \figcap{Cumulative recall as a function of retrieval depth $K$. The lexical structural pipeline accumulates evidence faster at low $K$ than the vector baseline.}
  \label{fig5_cumulative_recall}
\end{figure}

\subsection{Failure Mode Taxonomy}
We observe three failure classes. First, \textbf{Lexical Gap Refusal} occurs when user wording does not match canonical graph labels. Second, \textbf{Partial Path Failure} occurs when one anchor resolves and the other does not. Third, \textbf{Hub Saturation} appears when high degree institutions flood the candidate set, although this is strongly reduced by degree capping and stratified scoring.

\begin{figure}[htbp]
  \centering
  \safeincludegraphics[width=0.82\linewidth]{fig3_semantic_gap_implicit.png}
  \figcap{Lexical gap example. The semantic intent is valid, but the surface form differs from the canonical graph label, so anchor resolution fails and the system correctly refuses.}
  \label{fig3}
\end{figure}

\subsection{Token Efficiency Analysis}
Triplet centric contexts are much smaller than paragraph centric contexts. In the benchmark, GraphRAG prompts used roughly 65 to 75\% fewer evidence tokens than vector only baselines. This reduces both generation cost and positional degradation risk \cite{liu2024lostmiddle,openaipricing2026,anthropicpricing2026}.

\section{Discussion on Regulatory Compliance}
\subsection{Explainability in Analyst Workflows}
Every generated claim is accompanied by a \texttt{source\_chunk} pointer that identifies the exact filing passage supporting it. This enables compliance officers to inspect, challenge, and annotate evidence without requiring access to model internals. The provenance chain supports audit trails, model review, and challenge procedures \cite{ieeeeadd,kearns2019ethical}.

\subsection{Alignment with SR 11-7 and GDPR}
SR 11-7 requires conceptual soundness, ongoing monitoring, and effective challenge \cite{sr117}. Deterministic retrieval with explicit refusal improves all three properties because retrieval logic is interpretable, failure modes are observable, and every claim can be tested against its provenance chain. GDPR style accountability and explanation requirements are likewise better served when outputs are backed by inspectable source documents rather than opaque latent reasoning \cite{gdpr,nistairmf,euaiact2024}.

\subsection{High Risk Governance Alignment}
Automated financial analysis may fall under high risk governance expectations in emerging regulatory frameworks \cite{euaiact2024,occaiprinciples,fcaai2024}. The proposed architecture supports robustness through bounded path checks, supports accuracy through triplet level provenance, and supports human oversight through refusal and escalation.

\section{Threats to Validity}
The benchmark was authored in house, which may introduce selection bias. Results are reported on two issuers only, so external validity remains limited. Faithfulness evaluation involves human judgment, which introduces some subjectivity, although automated metrics such as RAGAS provide a useful complement \cite{ragas2023}. Finally, extraction uses a low temperature model configuration rather than strict symbolic parsing, so some run to run variance remains possible.

\section{Future Research}
The main limitation is lexical mismatch between user language and canonical node naming. Future work will add semantic entity resolution based on node embeddings and nearest neighbor candidate generation before deterministic traversal. Additional extensions include adaptive hop depth and cross filing temporal reasoning.

\section{Conclusion}
This paper presented a Lexical Structural GraphRAG architecture for SEC 10-K compliance analysis built on Triplet Atomicity and the Deterministic Path Requirement. The system replaces semantically retrieved prose chunks with typed provenance mapped facts and gates answer synthesis on the existence of a bounded evidence path in the knowledge graph. On a 35 query benchmark, the architecture achieves 100\% execution success, strong faithfulness, zero hallucination on unretrieved context, and materially smaller evidence context than vector only baselines. Its explicit failure modes, inspectable provenance chains, and governance aligned refusal policy make it well suited to high risk financial AI deployment.

\clearpage
\begin{thebibliography}{00}
\bibitem{basel}
Basel Committee on Banking Supervision, \textit{Basel III: A Global Regulatory Framework for More Resilient Banks and Banking Systems}, revised ed., 2011.

\bibitem{gdpr}
European Union, ``Regulation (EU) 2016/679 (General Data Protection Regulation),'' 2016.

\bibitem{sr117}
Board of Governors of the Federal Reserve System, \textit{Supervisory Guidance on Model Risk Management (SR 11-7)}, 2011.

\bibitem{bcbs239}
Basel Committee on Banking Supervision, \textit{Principles for Effective Risk Data Aggregation and Risk Reporting (BCBS 239)}, 2013.

\bibitem{iso27001}
ISO/IEC, ``ISO/IEC 27001:2022 Information Security, Cybersecurity and Privacy Protection,'' 2022.

\bibitem{pci40}
PCI Security Standards Council, ``PCI DSS v4.0,'' 2022.

\bibitem{lewis2020rag}
P. Lewis et al., ``Retrieval Augmented Generation for Knowledge Intensive NLP Tasks,'' in \textit{Advances in Neural Information Processing Systems}, 2020.

\bibitem{karpukhin2020dpr}
V. Karpukhin et al., ``Dense Passage Retrieval for Open Domain Question Answering,'' in \textit{Proceedings of EMNLP}, 2020.

\bibitem{colbert2020}
O. Khattab and M. Zaharia, ``ColBERT: Efficient and Effective Passage Search via Contextualized Late Interaction over BERT,'' in \textit{Proceedings of SIGIR}, 2020.

\bibitem{splade2021}
T. Formal et al., ``SPLADE v2: Sparse Lexical and Expansion Model for Information Retrieval,'' arXiv:2109.10086, 2021.

\bibitem{atlas2022}
G. Izacard et al., ``Few shot Learning with Retrieval Augmented Language Models,'' arXiv:2208.03299, 2022.

\bibitem{bender2021stochastic}
E. M. Bender et al., ``On the Dangers of Stochastic Parrots: Can Language Models Be Too Big?,'' in \textit{Proceedings of FAccT}, 2021.

\bibitem{sec10k}
U.S. Securities and Exchange Commission, ``Form 10-K,'' [Online]. Available: \url{https://www.sec.gov/forms}. [Accessed: 01-Apr-2026].

\bibitem{projectnotes}
M. A. Usmani, ``Project architectural notes and result analysis,'' internal manuscript, 2026.

\bibitem{hyde2023}
L. Gao et al., ``Precise Zero Shot Dense Retrieval without Relevance Labels,'' in \textit{Findings of ACL}, 2023.

\bibitem{edge2024graphrag}
D. Edge et al., ``From Local to Global: A GraphRAG Approach to Query Focused Summarization,'' arXiv:2404.16130, 2024.

\bibitem{attention2017}
A. Vaswani et al., ``Attention Is All You Need,'' in \textit{NeurIPS}, 2017.

\bibitem{bert2019}
J. Devlin et al., ``BERT: Pre training of Deep Bidirectional Transformers for Language Understanding,'' in \textit{NAACL HLT}, 2019.

\bibitem{roberta2019}
Y. Liu et al., ``RoBERTa: A Robustly Optimized BERT Pretraining Approach,'' arXiv:1907.11692, 2019.

\bibitem{sbert2019}
N. Reimers and I. Gurevych, ``Sentence BERT: Sentence Embeddings using Siamese BERT Networks,'' in \textit{EMNLP IJCNLP}, 2019.

\bibitem{t52020}
C. Raffel et al., ``Exploring the Limits of Transfer Learning with a Unified Text to Text Transformer,'' \textit{Journal of Machine Learning Research}, vol. 21, no. 140, 2020.

\bibitem{instructiontuning2022}
Y. Wei et al., ``Finetuned Language Models Are Zero Shot Learners,'' in \textit{ICLR}, 2022.

\bibitem{liu2024lostmiddle}
N. F. Liu et al., ``Lost in the Middle: How Language Models Use Long Contexts,'' \textit{Transactions of the Association for Computational Linguistics}, vol. 12, pp. 157--173, 2024.

\bibitem{truthfulqa2022}
S. Lin, J. Hilton, and O. Evans, ``TruthfulQA: Measuring How Models Mimic Human Falsehoods,'' in \textit{ACL}, 2022.

\bibitem{ragas2023}
S. Es et al., ``RAGAS: Automated Evaluation of Retrieval Augmented Generation,'' arXiv:2309.15217, 2023.

\bibitem{hotpotqa2018}
Z. Yang et al., ``HotpotQA: A Dataset for Diverse, Explainable Multi hop Question Answering,'' in \textit{EMNLP}, 2018.

\bibitem{fever2018}
J. Thorne et al., ``FEVER: A Large Scale Dataset for Fact Extraction and Verification,'' in \textit{NAACL HLT}, 2018.

\bibitem{edgargraph}
Y. Wang, J. Xu, and L. Zhang, ``Constructing Financial Knowledge Graphs from SEC Filings for Risk Analytics,'' in \textit{Proceedings of the IEEE International Conference on Big Data}, 2021.

\bibitem{finkgsurvey}
A. Hogan et al., ``Knowledge Graphs,'' \textit{ACM Computing Surveys}, vol. 54, no. 4, pp. 1--37, 2021.

\bibitem{opencyc2018}
A. M. F. et al., ``Knowledge Graphs in Financial Services: A Review,'' in \textit{IEEE Access}, 2020.

\bibitem{finbert2020}
D. Araci, ``FinBERT: Financial Sentiment Analysis with Pre trained Language Models,'' arXiv:1908.10063, 2019.

\bibitem{neo4jmanual}
Neo4j Inc., ``Neo4j Documentation,'' [Online]. Available: \url{https://neo4j.com/docs/}. [Accessed: 01-Apr-2026].

\bibitem{cypher2018}
F. Holzschuher and R. Peinl, ``Performance of Graph Query Languages: Comparison of Cypher, Gremlin and Native Access in Neo4j,'' in \textit{Datenbank Spektrum}, 2018.

\bibitem{gpt4report}
OpenAI, ``GPT-4 Technical Report,'' arXiv:2303.08774, 2023.

\bibitem{angles2018propertygraphs}
R. Angles et al., ``Foundations of Modern Query Languages for Graph Databases,'' \textit{ACM Computing Surveys}, vol. 50, no. 5, 2018.

\bibitem{msmarco2016}
N. Bajaj et al., ``MS MARCO: A Human Generated Machine Reading Comprehension Dataset,'' arXiv:1611.09268, 2016.

\bibitem{trec2020}
E. M. Voorhees and D. Harman, \textit{TREC: Experiment and Evaluation in Information Retrieval}. MIT Press, 2005.

\bibitem{robinson2015graphdb}
I. Robinson, J. Webber, and E. Eifrem, \textit{Graph Databases}, 2nd ed. O'Reilly Media, 2015.

\bibitem{monoT52020}
S. Nogueira et al., ``Document Ranking with a Pretrained Sequence to Sequence Model,'' in \textit{Findings of EMNLP}, 2020.

\bibitem{rankgpt2023}
T. Sun et al., ``Is ChatGPT Good at Search? Investigating Large Language Models as Re Ranking Agents,'' arXiv:2304.09542, 2023.

\bibitem{selfrag2023}
A. Asai et al., ``Self RAG: Learning to Retrieve, Generate, and Critique through Self Reflection,'' arXiv:2310.11511, 2023.

\bibitem{openaipricing2026}
OpenAI, ``API Pricing,'' [Online]. Available: \url{https://openai.com/api/pricing}. [Accessed: 01-Apr-2026].

\bibitem{anthropicpricing2026}
Anthropic, ``API Pricing,'' [Online]. Available: \url{https://www.anthropic.com/pricing}. [Accessed: 01-Apr-2026].

\bibitem{ieeeeadd}
IEEE, \textit{Ethically Aligned Design: A Vision for Prioritizing Human Well being with Autonomous and Intelligent Systems}, 1st ed., 2019.

\bibitem{kearns2019ethical}
M. Kearns and A. Roth, \textit{The Ethical Algorithm}. Oxford University Press, 2019.

\bibitem{nistairmf}
National Institute of Standards and Technology, \textit{AI Risk Management Framework (AI RMF 1.0)}, NIST AI 100-1, 2023.

\dots
\end{thebibliography}

\end{document}
